\documentclass[11pt]{article}

\usepackage[utf8]{inputenc}
\usepackage[T1]{fontenc}
\usepackage[margin=1in]{geometry}
\usepackage{booktabs}
\usepackage{array}
\usepackage{fancyvrb}
\usepackage[hidelinks]{hyperref}
\usepackage{url}

\newcolumntype{L}[1]{>{\raggedright\arraybackslash}p{#1}}

\title{gpu-mcp-server: Exposing GPU Telemetry to AI Agents\\
via the Model Context Protocol}

\author{Pavan Madduri\\
Independent Researcher\\
\texttt{https://orcid.org/0009-0007-3795-7593}}

\date{September 2026}

\begin{document}
\maketitle

\begin{abstract}
Large language model (LLM) agents are increasingly used to observe and operate
computing infrastructure, yet they have no standard, structured way to see the
state of the GPUs that dominate the cost and capacity of machine learning
systems. In practice an agent must invoke command-line tools such as
\texttt{nvidia-smi} and parse free-form text, which is brittle, inconsistent
across driver versions, and awkward to expose safely. We present
\texttt{gpu-mcp-server}, an open-source server that exposes NVIDIA GPU telemetry
to AI agents through the Model Context Protocol (MCP), an open standard for
connecting LLM applications to external tools. The server reads metrics directly
from the NVIDIA Management Library (NVML) and presents them as typed MCP tools
covering per-device utilization, memory, temperature, power, per-process usage,
and Multi-Instance GPU (MIG) partitions. It runs over stdio or streamable HTTP,
ships a container image and a Kubernetes deployment, and can be tested without a
physical GPU through a mock collector. By making GPU state a first-class,
structured tool rather than scraped command output, the server lets any
MCP-compatible agent reason about accelerator capacity, cost, and placement.
\end{abstract}

\section{Introduction}

GPUs are the scarce and expensive resource in modern machine learning
infrastructure, so accurate, real-time visibility into their utilization and
memory is essential for any system that reasons about capacity or cost. At the
same time, LLM-based agents are increasingly asked to monitor and operate
infrastructure. These two trends collide at a gap: agents have no standard way to
observe GPU state. The usual workaround---shelling out to \texttt{nvidia-smi} and
parsing its text output---is fragile across driver versions, hard to consume
programmatically, and difficult to expose to an agent with clear boundaries.

The Model Context Protocol (MCP) \cite{mcp} is an open protocol that standardizes
how LLM applications discover and call external tools through typed interfaces.
We use MCP to close the gap for GPUs. \texttt{gpu-mcp-server} \cite{gms}
implements the protocol and provides GPU telemetry as a small set of tools backed
by the NVIDIA Management Library (NVML) \cite{nvml}. Because it reads NVML rather
than scraping command output, the data is structured and stable, and any
MCP-compatible agent can query it without hardware-specific glue code.

\section{Background}

\paragraph{Model Context Protocol.} MCP defines a client--server protocol in
which a host application (for example, an assistant embedded in an editor or an
agent runtime) connects to servers that expose \emph{tools} with typed inputs and
outputs. This lets model-driven applications call capabilities in a uniform way
rather than through ad-hoc integrations.

\paragraph{NVIDIA Management Library.} NVML is the vendor library for querying and
managing NVIDIA GPU state, including utilization, memory, temperature, power, and
partition (MIG) information. \texttt{gpu-mcp-server} uses the official Go bindings
\cite{gonvml}, which load the library at runtime.

\section{Design}

\subsection{Tools}
The server registers a small set of MCP tools, summarized in
Table~\ref{tab:tools}. Each has a typed schema so agents receive structured
results rather than free-form text.

\begin{table}[h]
\centering
\begin{tabular}{L{0.30\textwidth} L{0.60\textwidth}}
\toprule
\textbf{Tool} & \textbf{Purpose} \\
\midrule
\texttt{list\_gpus} & Enumerate all GPUs with utilization and memory. \\
\texttt{get\_gpu\_metrics} & Detailed metrics for a specific GPU or MIG instance. \\
\texttt{gpu\_summary} & Aggregate statistics across all devices. \\
\texttt{get\_gpu\_processes} & Per-process (PID-level) GPU usage where supported. \\
\bottomrule
\end{tabular}
\caption{MCP tools exposed by the server.}
\label{tab:tools}
\end{table}

\subsection{Transports}
The server supports two transports: stdio, for local agent runtimes that launch
the server as a subprocess, and streamable HTTP, for networked or containerized
deployments. The same tool implementations serve both.

\subsection{Deployment}
The server is distributed as a container image and includes a Helm chart for
deployment on Kubernetes as a DaemonSet, so a cluster can expose per-node GPU
telemetry to agents. It also reports a build version for reproducibility.

\subsection{Testing without hardware}
Hardware access is abstracted behind a collector interface. A mock collector with
static device data lets the tool handlers be exercised on machines without a GPU,
so the server can be built and tested in ordinary CI.

\section{Use Cases}

Exposing GPU telemetry as MCP tools enables agent-driven workflows that are
otherwise awkward: an assistant can answer questions about live fleet
utilization, an operations agent can decide whether a node has spare capacity
before scheduling a job, and a cost-oriented agent can flag idle but allocated
GPUs. Because the server reads NVML directly, the same telemetry foundation
underlies the author's related Kubernetes GPU autoscaler \cite{kgs}, keeping
metric semantics consistent between agent-facing observability and
orchestrator-facing autoscaling.

\section{Related Work}

NVIDIA's \texttt{nvidia-smi} and DCGM expose GPU metrics, but as command-line
output or Prometheus time series rather than as agent-callable tools. General MCP
servers exist for many services, but not for direct GPU telemetry via NVML.
\texttt{gpu-mcp-server} is complementary to monitoring stacks: it does not replace
dashboards, but makes GPU state directly consumable by LLM agents through a
standard protocol.

\section{Conclusion}

AI agents increasingly operate infrastructure, but lacked a structured way to see
GPU state. By exposing NVML-backed telemetry as typed MCP tools over stdio and
HTTP, \texttt{gpu-mcp-server} gives agents reliable, real-time GPU visibility
without hardware-specific glue. The software is open source and available at
\cite{gms}.

\begin{thebibliography}{99}

\bibitem{mcp} Model Context Protocol. \url{https://modelcontextprotocol.io}

\bibitem{gms} P. Madduri. gpu-mcp-server: NVIDIA GPU metrics for AI agents over
the Model Context Protocol. Zenodo, 2026.
\url{https://doi.org/10.5281/zenodo.22866670}

\bibitem{nvml} NVIDIA Corporation. NVIDIA Management Library (NVML).
\url{https://developer.nvidia.com/management-library-nvml}

\bibitem{gonvml} NVIDIA Corporation. go-nvml: NVML bindings for Go.
\url{https://github.com/NVIDIA/go-nvml}

\bibitem{kgs} P. Madduri. keda-gpu-scaler: GPU-aware autoscaling for Kubernetes
from real hardware metrics. Zenodo, 2026.
\url{https://doi.org/10.5281/zenodo.22866676}

\end{thebibliography}

\end{document}
